\documentclass[11pt]{article}
\usepackage[margin=1in]{geometry}
\usepackage[T1]{fontenc}
\usepackage{lmodern}
\usepackage{microtype}
\usepackage{amsmath}
\usepackage{booktabs}
\usepackage{longtable}
\usepackage{array}
\usepackage{xcolor}
\usepackage{hyperref}

\hypersetup{colorlinks=true,linkcolor=blue!50!black,urlcolor=blue!50!black}
\setlength{\parindent}{0pt}
\setlength{\parskip}{0.45em}
\renewcommand{\arraystretch}{1.12}

\newcommand{\event}[1]{\subsection{#1}}
\newcommand{\prob}{\textbf{Firing probability.} }
\newcommand{\impact}{\textbf{Impact.} }
\newcommand{\eventmessage}{\textbf{Message.} }
\newcommand{\none}{No direct economic impact.}
\newcommand{\sched}[1]{\(\langle #1\rangle\)}
\newcommand{\impacttable}[5]{%
\begin{center}\small
\begin{tabular}{@{}l*{8}{r}@{}}
\toprule
Direct impact & Q0 & Q1 & Q2 & Q3 & Q4 & Q5 & Q6 & Q7\\
\midrule
Inflation & #1\\
Output gap & #2\\
Natural unemployment & #3\\
Equilibrium real rate & #4\\
Rate pressure & #5\\
\bottomrule
\end{tabular}
\end{center}}
\newcommand{\zeroimpact}{\impacttable{0&0&0&0&0&0&0&0}{0&0&0&0&0&0&0&0}{0&0&0&0&0&0&0&0}{0&0&0&0&0&0&0&0}{0&0&0&0&0&0&0&0}}

\title{Events in the PIRS Monetary-Policy Game}
\author{}
\date{September 2026}

\begin{document}
\maketitle
\tableofcontents

\section{How to read this document}

This document first describes the events and player actions currently in use.  Later sections propose future directions: inactive events that
could be enabled, new event candidates, and alternative event intensities and
news descriptions.  The existing Global Supply Shock is the benchmark for all
of the proposed intensity and description variations.

Every event below is presented in the same order: the message shown to the
player, one table containing all direct impacts, and finally a short probability
summary.  Table columns run from the event quarter (Q0) through seven later
quarters (Q7).  All values are additive percentage-point shocks. ``Rate pressure'' affects policy transmission ($R_{t}$) rather than directly changing the chosen policy rate. The output gap effects were defined in terms of unemployment in a previous version. Converting them to the current model lead to ``weird'' numbers like 1.43, instead of more round values.

Normally, at most one new random event is announced in a quarter, while effects
from earlier events continue.  The approximate probabilities below assume the
standard probability scale.  Difficulty imposes a general gap between random
events: 20 quarters in Principles, 10 in Senior, and none in Central Banker mode.

\section{Random and informational events}

\event{Financial Crisis}
\eventmessage ``A financial crisis leads to a credit crunch and reduced investment.''
\impacttable{-0.2&-0.2&0&0&0&0&0&0}{-0.14&-1.43&-1.43&-0.71&-0.71&0&0&0}{0.2&0.2&0.5&0&0&-0.2&-0.2&-0.5}{-2&-1&1&0.5&0.5&0.5&0.5&0}{0&0&0&0&0&0&0&0}
\prob Approximately 1 percent in normal conditions.  Recent rate increases and a run of rates below 1 percent raise the chance; an occurrence in the last four quarters prevents a repeat.

\event{Major Financial Crisis}
\eventmessage ``Panic spreads through global markets, with commentators evoking the catastrophic collapse of 1929 as systemic crisis looms.''
\impacttable{-0.2&-0.5&-0.3&0&0&0&0&0}{-1.43&-2.86&-4.29&-5.71&-4.29&-3.57&-2.86&-1.43}{0&1&1&0&-0.25&-0.25&-0.5&-1}{-5&-2&1&1&1&1&1&2}{0&0&0&0&0&0&0&0}
\prob Approximately 0.25 percent in normal conditions.  A recent Financial Crisis sharply raises the chance, especially when the policy rate is at least as high as inflation; active quantitative easing halves this escalation risk.

\event{High Trust}
\eventmessage ``Public confidence in the Central Bank's commitment to low inflation is strong. There is room to maneuver without losing credibility.''
\zeroimpact
\prob Approximately 20 percent when reputation is above 0.8.  It cannot repeat within eight quarters.

\event{Moderate Trust}
\eventmessage ``Confidence in the Central Bank is moderate and markets are cautions.Consistent policy could strengthen credibility.''
\zeroimpact
\prob Approximately 90 percent when reputation is between 0.4 and 0.7.  It cannot repeat within eight quarters.

\event{Low Trust}
\eventmessage ``Public confidence in the Central Bank is weak. Specialists warn credibility must be rebuilt through clear policy signals.''
\zeroimpact
\prob Approximately 20 percent when reputation is below 0.3.  It cannot repeat within eight quarters.

\event{Pressure for Lower Interest Rates}
\eventmessage ``Public and markets demand the Central Bank to cut rates to stimulate the economy.''
\zeroimpact
\prob Approximately 10 percent when the policy rate is above 2 percent and unemployment is above 7 percent.  It cannot repeat within eight quarters.

\event{High Inflation Warning}
\eventmessage ``Inflation has crossed 10\%. Analysts fear a spiral unless Central Bank restores stability.''
\zeroimpact
\prob Approximately 20 percent when inflation is rising and is between 10 and 20 percent.  It cannot repeat within eight quarters.

\event{Hyperinflation Risk}
\eventmessage ``Inflation has crossed 100\%! Analysts fear hyperinflation while massive protests take the streets.''
\zeroimpact
\prob Approximately 10 percent when inflation exceeds 100 percent and has recently been rising.  It cannot repeat within eight quarters.

\event{Technological Boom}
\eventmessage ``A dramatic acceleration in technological progress reshapes the economy, echoing past transformations but with modern global scale.''
\impacttable{0.1&0.2&0&-0.5&-1&-1&-0.5&0}{0.43&0.29&-0.71&-0.71&-0.71&-0.71&-0.14&-0.14}{-0.3&-0.3&0.5&0.5&1&-0.1&-0.1&-0.1}{0.1&0.2&0&0&-0.2&-0.1&0&0}{0&0&0&0&0&0&0&0}
\prob Approximately 1 percent in normal conditions.

\event{Pandemic Outbreak}
\eventmessage ``A public health crisis disrupts supply chains and labor markets''
\impacttable{1&2&2&0.5&0.5&0.5&0.2&0}{-10&-8.57&10&2.86&1.43&0&0&0}{2&1&-1&-1&-1&-1&0&0}{0&0&0&0&0&0&0&0}{0&0&0&0&0&0&0&0}
\prob Approximately 0.5 percent in normal conditions.

\event{Natural Disaster}
\eventmessage ``Severe natural events damage infrastructure and disrupt production.''
\impacttable{1&2&0&0&0&0&0&0}{-0.43&-0.14&0&0&0&0&0&0}{0.2&0.1&0&0&0&0&0&-0.3}{0&0&0&0&0&0&0&0}{0&0&0&0&0&0&0&0}
\prob Approximately 1 percent in normal conditions.  It cannot repeat within eight quarters.

\event{Global Supply Shock}
\eventmessage ``Geopolitical conflict disrupts energy supply, pushing up costs across the economy.''
\impacttable{1.5&2.5&1.5&0.5&0&0&0&0}{-0.43&-0.71&-0.43&-0.14&0&0&0&0}{1&3&3&0&-1&-1&-1&-1}{0.2&0.2&0.1&-0.1&-0.2&-0.2&0&0}{0&0&0&0&0&0&0&0}
\prob Approximately 0.8 percent in normal conditions.

\event{Fiscal Deficit}
\eventmessage ``Deficits support short-term demand, but raise fears of unsustainable debt and inflation.''
\impacttable{0.2&0.1&0.1&0&0&0&0&0}{0.71&1.14&0.71&0&0&0&0&0}{0&0&0&0&0&0&0&0}{1&-0.1&-0.1&-0.1&-0.1&-0.1&-0.2&-0.3}{0&0&0&0&0&0&0&0}
\prob Approximately 2 percent in normal conditions, rising to about 12 percent when inflation is below 1 percent and unemployment exceeds 10 percent.  A recent occurrence usually prevents a repeat; Principles mode may also trigger it to fight deflation and excess unemployment.

\event{Spending Wave}
\eventmessage ``Massive public spending is hailed as a path to jobs and prosperity, though critics warn of hyperinflation.''
\impacttable{0.7&1&1&0.5&0.3&0&0&0}{2.86&2.86&2.86&1.43&1.43&0&0&0}{0&0&0&0&0&0&0&0}{4&0&0&0&-1&-1&-1&-1}{0&0&0&0&0&0&0&0}
\prob Approximately 0.5 percent in normal conditions, rising to about 10.5 percent during severe deflation and unemployment.  A recent occurrence strongly suppresses another; Principles mode may also trigger it to fight deflation and excess unemployment.

\event{Fiscal Surplus}
\eventmessage ``Balanced books inspire confidence in sustainability, even as critics caution against tightening too far.''
\impacttable{0&0&0&0&0&0&0&0}{-0.71&-0.29&0&0&0&0&0&0}{0.5&0&0&0&0&0&0&0}{-1&0.1&0.1&0.1&0.1&0.1&0.2&0.3}{0&0&0&0&0&0&0&0}
\prob Approximately 1 percent in normal conditions.  It cannot repeat within eight quarters; Principles mode may also trigger it when inflation exceeds 10 percent.

\event{Research on Policy Lags}
\eventmessage ``New research suggests real interest rate changes take about 1 year to impact employment and about 2 years to impact inflation.''
\zeroimpact
\prob Expected to appear during user quarters 41--44 in Senior or Central Banker difficulty.  It cannot repeat within eight quarters and may be delayed by the general event cooldown.

\event{Hawk Economist Calls for Rate Hike}
\eventmessage ``'It is time for a decisive hike: only by holding rates well above inflation can we discipline the labour market, restore the bank's authority, and crush inflation.', she says''
\zeroimpact
\prob Approximately 10 percent for every percentage point that inflation is above 3 percent, up to 60 percent.  A recent occurrence prevents a repeat, and the event does not fire when the policy rate is already more than 3 points above inflation.

\section{Player-triggered policy events}

These actions occur when the player selects them and the availability rules pass.  They are available only in Central Banker mode, and only one player action may be initiated in a quarter.

\event{Quantitative Easing}
\eventmessage Headline: ``Central Bank Launches Asset Purchases.''  Detail: ``The Central Bank has announced a new asset-purchase programme to support demand and ease financial conditions.''
\impacttable{0.05&0.1&0.065&0.04&0.025&0.015&0.01&0.005}{0.5&1&0.65&0.4&0.25&0.15&0.1&0.05}{0&0&0&0&0&0&0&0}{0&0&0&0&0&0&0&0}{0&0&0&0&0&0&0&0}
\prob Not random: it fires when selected.  It has no action-specific cooldown, and while active it reduces escalation from a Financial Crisis to a Major Financial Crisis.

\event{Announce Future High Rates}
\eventmessage Headline: ``Central Bank Signals Higher Rates Ahead.''  Detail: ``The Central Bank has signalled that interest rates are likely to remain higher in the coming quarters.''
\impacttable{0&0&0&0&0&0&0&0}{0&0&0&0&0&0&0&0}{0&0&0&0&0&0&0&0}{0&0&0&0&0&0&0&0}{1&1&1&1&0&0&0&0}
\prob Not random: it fires when selected, provided reputation is above 0.70.  The same announcement has a four-quarter cooldown.

\event{Announce Future Low Rates}
\eventmessage Headline: ``Central Bank Signals Lower Rates Ahead.''  Detail: ``The Central Bank has signalled that interest rates are likely to remain lower in the coming quarters.''
\impacttable{0&0&0&0&0&0&0&0}{0&0&0&0&0&0&0&0}{0&0&0&0&0&0&0&0}{0&0&0&0&0&0&0&0}{-1&-1&-1&-1&0&0&0&0}
\prob Not random: it fires when selected.  The same announcement has a four-quarter cooldown.


\section{Currently inactive explainers}

These events are currently inactive but could be incorporated as part of a "tutorial" option and/or at the principles level.

\event{Explainer: Unemployment and Inflation}
\eventmessage ``Lesson: Unemployment measures the share of people actively looking for work who cannot find jobs. When unemployment is very low, wages and demand can rise faster, adding inflation pressure. When it is high, inflation usually cools because spending power weakens.''
\zeroimpact
\prob 0 percent: this explainer is currently disabled.

\event{Explainer: Real Interest Rates and Employment}
\eventmessage ``Lesson: The real interest rate is approximately the policy rate minus inflation. Higher real rates make borrowing costlier and tend to reduce aggregate demand, which can raise unemployment. Lower real rates usually support demand and employment, with effects that arrive gradually.''
\zeroimpact
\prob 0 percent: this explainer is currently disabled.

\event{Explainer: Trust and Expectations Anchoring}
\eventmessage ``Lesson: If people trust the central bank, inflation expectations stay anchored near target. Anchored expectations reduce self-fulfilling inflation spirals and make policy less costly for jobs.''
\zeroimpact
\prob 0 percent: this explainer is currently disabled.

\event{Explainer: Natural Unemployment Rate}
\eventmessage ``Lesson: The natural unemployment rate is the level consistent with stable inflation in the medium run, reflecting frictions like job search and skill mismatch. Policy can move unemployment around it temporarily, but structural reforms influence it more durably.''
\zeroimpact
\prob 0 percent: this explainer is currently disabled.

\event{Explainer: Nominal vs Real Variables}
\eventmessage ``Lesson: Nominal variables are measured in current prices, while real variables adjust for inflation. Policy decisions based only on nominal rates can be misleading when inflation shifts quickly.''
\zeroimpact
\prob 0 percent: this explainer is currently disabled.

\section{New Events Proposed (by AI)}

The probabilities in this section are proposed per-quarter hazards, not targets
chosen to balance positive and negative shocks or to fill an event quota.  They
are calibrated to the rough real-world frequency of the episode described.
Conditional adjustments use only state that the game already records:
inflation, unemployment, policy and equilibrium rates, reputation, their
histories, and past events.


\event{Housing Market Correction }
\eventmessage ``House prices fall sharply as stretched borrowers pull back and residential construction slows.''
\impacttable{-0.1&-0.2&-0.2&0&0&0&0&0}{-0.4&-0.7&-0.5&-0.2&0&0&0&0}{0&0&0&0&0&0&0&0}{-0.3&-0.2&0&0.1&0.1&0&0&0}{0&0&0&0&0&0&0&0}
\prob Approximately 0.3 percent per quarter in normal conditions.  A sequence of policy-rate increases or a policy rate well above its recent average could raise the chance; it should not repeat within twelve quarters.

\event{Consumer Confidence Surge }
\eventmessage ``Households turn more optimistic about their finances, lifting retail spending and demand for big-ticket goods.''
\impacttable{0.1&0.2&0.2&0.1&0&0&0&0}{0.4&0.6&0.4&0.2&0&0&0&0}{0&0&0&0&0&0&0&0}{0.1&0.1&0&0&0&0&0&0}{0&0&0&0&0&0&0&0}
\prob Approximately 1 percent per quarter in normal conditions.  Falling unemployment over the previous four quarters and inflation moving toward target could each raise the chance modestly; it should not repeat within eight quarters.

\event{Productivity Slowdown }
\eventmessage ``Businesses report weaker productivity growth as investment projects disappoint and operating costs rise.''
\impacttable{0.3&0.5&0.4&0.3&0.2&0.1&0&0}{-0.2&-0.4&-0.5&-0.4&-0.2&0&0&0}{0.1&0.2&0.2&0.1&0&0&0&0}{-0.2&-0.2&-0.1&0&0&0&0&0}{0&0&0&0&0&0&0&0}
\prob Approximately 0.6 percent per quarter in normal conditions.  Because the game does not track productivity or investment, this should remain an unconditional draw apart from a twelve-quarter repeat restriction.

\event{Export Boom }
\eventmessage ``Foreign orders jump, giving exporters a surprise boost and increasing demand for domestic workers and equipment.''
\impacttable{0.1&0.3&0.3&0.2&0.1&0&0&0}{0.5&0.8&0.7&0.4&0.2&0&0&0}{0&0&0&0&0&0&0&0}{0.2&0.2&0.1&0&0&0&0&0}{0&0&0&0&0&0&0&0}
\prob Approximately 0.8 percent per quarter in normal conditions.  Because the game does not track exports or foreign demand, this should remain an unconditional draw apart from an eight-quarter repeat restriction.

\event{Banking Cyberattack }
\eventmessage ``A major cyberattack disrupts bank payments and briefly tightens credit while financial institutions restore their systems.''
\impacttable{0&-0.1&0&0&0&0&0&0}{-0.6&-0.5&-0.2&0&0&0&0&0}{0&0&0&0&0&0&0&0}{-0.4&-0.2&0&0&0&0&0&0}{0&0&0&0&0&0&0&0}
\prob Approximately 0.1 percent per quarter in normal conditions.  This system-wide disruption should be an unconditional, exceptionally rare event and should not repeat within twenty quarters.

\event{Immigration Wave }
\eventmessage ``A rise in immigration expands the workforce and consumer base, easing some hiring shortages while increasing near-term demand.''
\impacttable{0.1&0.1&0&-0.1&-0.1&0&0&0}{0.2&0.4&0.5&0.4&0.2&0&0&0}{-0.2&-0.3&-0.3&-0.2&-0.1&0&0&0}{0.1&0.1&0.1&0&0&0&0&0}{0&0&0&0&0&0&0&0}
\prob Approximately 0.5 percent per quarter in normal conditions.  Unemployment below 4 percent could raise the chance modestly, using labor-market tightness as the game's observable proxy; it should not repeat within twelve quarters.

\event{Nationwide Wage Agreement }
\eventmessage ``Employers and unions reach a broad wage agreement, supporting household incomes but adding to firms' labor costs.''
\impacttable{0.3&0.5&0.4&0.2&0.1&0&0&0}{0.2&0.3&0.2&0&0&0&0&0}{0&0&0&0&0&0&0&0}{0.1&0.1&0&0&0&0&0&0}{0&0&0&0&0&0&0&0}
\prob Approximately 0.6 percent per quarter in normal conditions.  Inflation above 4 percent or unemployment below 4 percent could raise the chance, while the combination could raise it further; it should not repeat within twelve quarters.

\event{Crop Harvest Windfall }
\eventmessage ``An unusually strong harvest lowers food prices and raises rural incomes, offering households temporary relief.''
\impacttable{-0.5&-0.7&-0.4&-0.2&0&0&0&0}{0.1&0.2&0.2&0.1&0&0&0&0}{0&0&0&0&0&0&0&0}{0&0&0&0&0&0&0&0}{0&0&0&0&0&0&0&0}
\prob Approximately 1 percent per quarter in normal conditions.  Weather and harvest conditions are not tracked, so this should be an unconditional draw, independent of monetary policy, with an eight-quarter repeat restriction.

\event{Corporate Investment Wave }
\eventmessage ``Businesses announce a wave of factory, equipment, and software investment, strengthening demand before new capacity comes online.''
\impacttable{0.1&0.2&0.2&0&-0.2&-0.3&-0.2&0}{0.4&0.7&0.8&0.6&0.4&0.2&0.1&0}{0&0&0&0&-0.1&-0.1&-0.1&0}{0.3&0.4&0.3&0.2&0.1&0&0&0}{0&0&0&0&0&0&0&0}
\prob Approximately 0.8 percent per quarter in normal conditions.  Inflation within one percentage point of target and a policy rate no more than two points above inflation plus the equilibrium real rate could each raise the chance modestly; it should not repeat within ten quarters.

\event{Sovereign Debt Scare }
\eventmessage ``Investors demand higher returns on government debt as concerns about the public finances spill into private borrowing markets.''
\impacttable{0.2&0.3&0.2&0.1&0&0&0&0}{-0.3&-0.6&-0.5&-0.3&-0.1&0&0&0}{0&0.1&0.1&0&0&0&0&0}{0.8&0.6&0.4&0.2&0.1&0&0&0}{0&0&0&0&0&0&0&0}
\prob Approximately 0.2 percent per quarter in normal conditions.  A Fiscal Deficit or Spending Wave in the previous eight quarters could raise the chance, using only the recorded event history; it should not repeat within sixteen quarters.

\section{Varying current events}

This section uses the existing Global Supply Shock as a test case for the other
possible interpretations of the feedback.  These are alternative design options,
not additions that should all be active together.  Cosmetic and explanatory
variants retain the current event's impacts and probability.  Mechanical variants
would instead create distinct intensities.

\subsection{Cosmetic variations}

These messages preserve the present length, tone, mechanics, and approximate
0.8-percent probability.  One could be selected at random whenever the existing
Global Supply Shock fires.



\subsubsection{Current version}
\eventmessage ``Geopolitical conflict disrupts energy supply, pushing up costs across the economy.''

\subsubsection{Cosmetic variation 1}
\eventmessage ``International tensions interrupt fuel shipments, raising transport and production costs throughout the economy.''

\subsubsection{Cosmetic variation 2}
\eventmessage ``A disruption to global energy trade leaves businesses facing shortages and a sudden increase in input prices.''

\subsection{Mechanical variations}

These alternatives give supply shocks different intensities.  Their names and
messages make the size visible rather than asking the player to infer it from the
charts.  A shared eight-quarter restriction across all supply-shock intensities
would prevent them from appearing as repetitive back-to-back stories.

\subsection{Current Version}

\impacttable{1.5&2.5&1.5&0.5&0&0&0&0}{-0.43&-0.71&-0.43&-0.14&0&0&0&0}{1&3&3&0&-1&-1&-1&-1}{0.2&0.2&0.1&-0.1&-0.2&-0.2&0&0}{0&0&0&0&0&0&0&0}
\prob Approximately 0.8 percent in normal conditions.

\subsubsection{Limited Supply Disruption}
\eventmessage ``Temporary shipping and energy disruptions lift costs, but markets expect the shortages to remain contained.''
\impacttable{0.6&1&0.6&0.2&0&0&0&0}{-0.14&-0.29&-0.14&0&0&0&0&0}{0.3&0.8&0.8&0&-0.3&-0.3&-0.3&-0.2}{0.1&0.1&0&-0.1&-0.1&0&0&0}{0&0&0&0&0&0&0&0}
\prob Approximately 1 percent in normal conditions.  This is slightly more common than the current event because its effects are weaker and shorter-lived.

\subsubsection{Severe Global Supply Crisis}
\eventmessage ``A widening geopolitical crisis cuts vital energy supplies, forcing factories to reduce production as prices surge worldwide.''
\impacttable{2.5&4&3&1.5&0.8&0.3&0&0}{-0.71&-1.14&-0.86&-0.43&-0.14&0&0&0}{1.5&4&4&1&-1&-1.5&-1.5&-1.5}{0.4&0.4&0.2&-0.2&-0.3&-0.2&0&0}{0&0&0&0&0&0&0&0}
\prob Approximately 0.2 percent in normal conditions.  A recent ordinary supply shock or elevated geopolitical stress could raise the chance, while the shared repeat restriction prevents rapid recurrence.

\subsection{Longer and more colorful variations}

These alternatives retain the current mechanics and probability but make the
news more vivid.  The first is approximately twice the current message's length
and the second approximately four times its length.  They should be used
cautiously: their extra atmosphere may compete with the policy information the
player needs.

\subsubsection{Colorful variation 1: approximately 2x longer}
\eventmessage ``Geopolitical conflict strands tankers offshore and forces factories to ration fuel, lifting transport, heating, and production costs throughout the economy.''

\subsubsection{Colorful variation 2: approximately 4x longer}
\eventmessage ``Tankers idle outside shuttered ports as a widening geopolitical conflict severs major energy routes. Fuel wholesalers scramble for scarce cargoes, factories shorten shifts, and freight operators add emergency surcharges. With heating, transport, and production bills all climbing, businesses warn that customers will soon face higher prices.''

\subsection{More pedagogical variations}

These alternatives also retain the current mechanics and probability.  Each
briefly connects the story to aggregate supply and to the policy trade-off.

\subsubsection{Pedagogical variation 1}
\eventmessage ``A disruption to global energy supplies raises firms' production costs. This shifts short-run aggregate supply inward, tending to increase inflation while reducing output and employment.''

\subsubsection{Pedagogical variation 2}
\eventmessage ``Imported energy has become scarcer and more expensive. The economy therefore faces a negative supply shock: prices rise even as activity weakens, leaving monetary policy with a harder inflation--employment trade-off.''


\end{document}
